\documentclass[11pt,letterpaper]{article}
\usepackage{../../Common/course_notes}
\usepackage{../../Common/course_macros}

\begin{document}
\setlength{\abovedisplayskip}{6pt}
\setlength{\belowdisplayskip}{6pt}
\setlength{\textfloatsep}{10pt}
\setlength{\intextsep}{8pt}
\color{blackmain}

\begin{center}
{\LARGE\textbf{Geometría y Física I: Mecánica Clásica y Formas de Cartan}}\\[0.25cm]
{\Large\textcolor{orangeacc}{\textbf{Tarea 2: Construcción Geométrica del Arsenal de Índices}}}\\[0.2cm]
{\large\textcolor{graymuted}{Nivel: Licenciatura en Física --- Módulo 0: Arsenal Indicial}}\\[0.25cm]
\end{center}

\vspace{0.2cm}

\noindent\textbf{Instrucciones:} Resuelve detalladamente los siguientes problemas. El objetivo de esta tarea no es aplicar fórmulas a ciegas, sino \textbf{construir y deducir} las propiedades algebraicas a partir de la geometría de volúmenes, áreas y simetrías. Utiliza en todo momento el convenio de suma de Einstein.

% =====================================================================
\manualsection{Parte A: Construcción Geométrica en 2D y 3D}

\begin{infobox}{Problema 1: Construyendo el Símbolo de Levi-Civita en 2D (Áreas Orientadas)}
En clase construimos $\epsilon_{ijk}$ en $\R^3$ como el volumen orientado de tres vectores base. Ahora construiremos su análogo en el plano $\R^2$, donde los índices solo toman dos valores: $i, j \in \{1, 2\}$.

Definimos $\epsilon_{ij}$ como el \textbf{área con signo} del paralelogramo formado por los vectores base $(\hat{e}_i, \hat{e}_j)$ del plano (positivo si el giro de $\hat{e}_i$ a $\hat{e}_j$ es antihorario, negativo si es horario, y cero si el paralelogramo se aplasta).
\begin{enumerate}[label=(\alph*)]
    \item Usando esta definición geométrica, encuentra las $2^2 = 4$ componentes de $\epsilon_{ij}$ y escríbelas como una matriz $2 \times 2$:
    \[
    (\epsilon_{ij}) = \begin{pmatrix} \epsilon_{11} & \epsilon_{12} \\ \epsilon_{21} & \epsilon_{22} \end{pmatrix}.
    \]
    \item Dados dos vectores en el plano $\vec{u} = u_i\hat{e}_i$ y $\vec{v} = v_j\hat{e}_j$, expande la suma doble $\text{Área}(\vec{u}, \vec{v}) = \epsilon_{ij} u_i v_j$ y verifica que reproduce el determinante $2 \times 2$ clásico $u_1 v_2 - u_2 v_1$.
    \item \textbf{Deducción por conteo de casos:} Analiza el producto $\epsilon_{ij}\epsilon_{mn}$ en $\R^2$ (sin índices sumados). Para que este producto no sea cero, tanto $(i,j)$ como $(m,n)$ deben ser permutaciones de $(1,2)$. Siguiendo el mismo razonamiento de casos visto en clase, demuestra que:
    \[
    \epsilon_{ij}\epsilon_{mn} = \delta_{im}\delta_{jn} - \delta_{in}\delta_{jm}.
    \]
    \item Contrae el primer índice haciendo $m = i$ en la identidad del inciso (c) (recuerda que en $\R^2$ se cumple $\delta_{ii} = \delta_{11} + \delta_{22} = 2$) y demuestra que $\epsilon_{ij}\epsilon_{in} = \delta_{jn}$.
\end{enumerate}
\end{infobox}

\begin{infobox}{Problema 2: Anatomía del Volumen Elemental $\epsilon_{ijk}$ en $\R^3$}
En $\R^3$, construimos $\epsilon_{ijk} \equiv \hat{e}_i \cdot (\hat{e}_j \times \hat{e}_k)$, el cual posee $3^3 = 27$ componentes.
\begin{enumerate}[label=(\alph*)]
    \item Escribe explícitamente las tres matrices $3 \times 3$ obtenidas al fijar el primer índice en $i = 1, 2, 3$:
    \[
    M^{(1)}_{jk} = \epsilon_{1jk}, \qquad M^{(2)}_{jk} = \epsilon_{2jk}, \qquad M^{(3)}_{jk} = \epsilon_{3jk}.
    \]
    \item Observa que las tres matrices son antisimétricas ($M^{(i)}_{jk} = -M^{(i)}_{kj}$). Explica geométricamente por qué la diagonal principal de todas ellas está compuesta exclusivamente por ceros.
\end{enumerate}
\end{infobox}

% =====================================================================
\manualsection{Parte B: Escalera de Contracciones y Demolición de Identidades}

\begin{infobox}{Problema 3: La Escalera de Contracciones de Levi-Civita}
Partiendo de la identidad maestra que construimos en clase para un índice contraído en $\R^3$:
\[
\epsilon_{ijk}\epsilon_{imn} = \delta_{jm}\delta_{kn} - \delta_{jn}\delta_{km},
\]
deduce qué ocurre cuando contraemos más índices:
\begin{enumerate}[label=(\alph*)]
    \item \textbf{Dos índices contraídos:} Haz $m = j$ (es decir, calcula $\epsilon_{ijk}\epsilon_{ijn}$). Usa el hecho de que en $\R^3$ la traza es $\delta_{jj} = 3$ y que $\delta_{jn}\delta_{kj} = \delta_{kn}$ para demostrar que:
    \[
    \mathbf{\epsilon_{ijk}\epsilon_{ijn} = 2\delta_{kn}}.
    \]
    Explica combinatoriamente de dónde sale ese factor $2$: si fijas el último índice $k=n$ (por ejemplo $k=n=3$), ¿cuántas parejas $(i,j)$ no nulas existen?
    \item \textbf{Tres índices contraídos:} Ahora haz $n = k$ en el resultado anterior para calcular la suma total $\epsilon_{ijk}\epsilon_{ijk}$. ¿Por qué el resultado coincide con el factorial $3! = 6$?
\end{enumerate}
\end{infobox}

\begin{infobox}{Problema 4: Producto Escalar de dos Productos Cruz e Identidad de Lagrange}
\begin{enumerate}[label=(\alph*)]
    \item Dados cuatro vectores $\vec{A}, \vec{B}, \vec{C}, \vec{D} \in \R^3$, escribe el producto escalar $(\vec{A}\times\vec{B})\cdot(\vec{C}\times\vec{D})$ como $(\vec{A}\times\vec{B})_i (\vec{C}\times\vec{D})_i$ y demuestra en dos renglones que:
    \[
    (\vec{A} \times \vec{B}) \cdot (\vec{C} \times \vec{D}) = (\vec{A}\cdot\vec{C})(\vec{B}\cdot\vec{D}) - (\vec{A}\cdot\vec{D})(\vec{B}\cdot\vec{C}).
    \]
    \item Evalúa el caso particular $\vec{C} = \vec{A}$ y $\vec{D} = \vec{B}$ para obtener la \textbf{Identidad de Lagrange}:
    \[
    |\vec{A} \times \vec{B}|^2 = |\vec{A}|^2|\vec{B}|^2 - (\vec{A}\cdot\vec{B})^2.
    \]
    Sustituye $\vec{A}\cdot\vec{B} = |\vec{A}||\vec{B}|\cos\theta$ y deduce geométricamente la fórmula $|\vec{A}\times\vec{B}| = |\vec{A}||\vec{B}||\sin\theta|$.
\end{enumerate}
\end{infobox}

\begin{infobox}{Problema 5: ¿Es asociativo el Producto Cruz? La Identidad de Jacobi}
A diferencia de la multiplicación de números reales, el producto cruz \textbf{no es asociativo}: en general $\vec{A}\times(\vec{B}\times\vec{C}) \neq (\vec{A}\times\vec{B})\times\vec{C}$. En su lugar, satisface una ley cíclica fundamental llamada \textbf{Identidad de Jacobi}.

Usando la fórmula BAC--CAB demostrada en clase, expande cada uno de los tres términos siguientes y demuestra en tres renglones que su suma se anula idénticamente:
\[
\vec{A}\times(\vec{B}\times\vec{C}) + \vec{B}\times(\vec{C}\times\vec{A}) + \vec{C}\times(\vec{A}\times\vec{B}) = \vec{0}.
\]
\end{infobox}

% =====================================================================
\manualsection{Parte C: Física y Cálculo Diferencial con Índices}

\begin{infobox}{Problema 6: Campos Centrales y Conservación del Momento Angular}
El momento angular de una partícula es $\vec{L} = \vec{r} \times \vec{p}$, cuya componente $i$-ésima es $L_i = \epsilon_{ijk} x_j p_k$.
\begin{enumerate}[label=(\alph*)]
    \item Deriva respecto al tiempo usando la regla de Leibniz: $\dot{L}_i = \epsilon_{ijk} \dot{x}_j p_k + \epsilon_{ijk} x_j \dot{p}_k$.
    \item Como $\vec{p} = m\dot{\vec{r}}$, el primer término es $m\,\epsilon_{ijk} \dot{x}_j \dot{x}_k$. Demuestra mediante el principio de contracción simétrico-antisimétrico que este término es idénticamente cero.
    \item Por la segunda ley de Newton, $\dot{p}_k = F_k$. Si la partícula se mueve bajo una fuerza gravitatoria o electrostática central $\vec{F} = f(r)\vec{r}$ (donde $F_k = f(r)x_k$), demuestra en un renglón que $\dot{L}_i = 0$.
\end{enumerate}
\end{infobox}

\begin{infobox}{Problema 7: El Operador Derivada $\partial_i$ y el Precursor de $d^2 = 0$}
Denotemos la derivada parcial respecto a la coordenada $x_i$ como $\partial_i \equiv \frac{\partial}{\partial x_i}$. En cálculo de varias variables, el \textbf{Teorema de Schwarz} asegura que las derivadas cruzadas de una función suave conmutan:
\[
\partial_i \partial_j f = \partial_j \partial_i f \quad (\text{el operador } \partial_i \partial_j \text{ es \textbf{simétrico} en } i, j).
\]
En notación de índices, el gradiente es $(\vec{\nabla}f)_i = \partial_i f$, la divergencia es $\vec{\nabla}\cdot\vec{V} = \partial_i V_i$ y el rotacional es $(\vec{\nabla}\times\vec{V})_i = \epsilon_{ijk} \partial_j V_k$.
\begin{enumerate}[label=(\alph*)]
    \item \textbf{El rotacional de un gradiente:} Escribe la componente $i$-ésima de $\vec{\nabla} \times (\vec{\nabla}f)$ en índices y demuestra en un solo renglón, usando la simetría de Schwarz frente a la antisimetría de $\epsilon_{ijk}$, que:
    \[
    [\vec{\nabla} \times (\vec{\nabla}f)]_i = \epsilon_{ijk} \partial_j \partial_k f = 0.
    \]
    \item \textbf{La divergencia de un rotacional:} Escribe $\vec{\nabla} \cdot (\vec{\nabla} \times \vec{V})$ en índices y demuestra por el mismo argumento que:
    \[
    \vec{\nabla} \cdot (\vec{\nabla} \times \vec{V}) = \epsilon_{ijk} \partial_i \partial_j V_k = 0.
    \]
\end{enumerate}
\end{infobox}

\end{document}
